\documentclass[11pt]{article}

\usepackage[margin=1in]{geometry}
\usepackage[utf8]{inputenc}
\usepackage[T1]{fontenc}
\usepackage{hyperref}
\usepackage{url}
\usepackage{booktabs}
\usepackage{amsmath,amssymb,amsfonts}
\usepackage{algorithm}
\usepackage{algorithmic}
\usepackage{graphicx}
\usepackage{microtype}
\usepackage{xcolor}
\usepackage{subcaption}
\usepackage{natbib}
\usepackage{multirow}

\title{TurboQuant on Quantized Models:\\Solving Compound Quantization Error in GGUF LLMs}

\author{
  Onur G\"oky{\i}ld{\i}z \\
  BHI Research \\
  \texttt{onur@bhi.dev} \\
  \url{https://github.com/onur-gokyildiz-bhi/tq-kv}
}

\begin{document}
\maketitle

\begin{abstract}
TurboQuant~\citep{zandieh2026turboquant} achieves near information-theoretic optimal KV cache compression through randomized Hadamard rotation and Lloyd-Max scalar quantization, delivering up to 15$\times$ compression with 0.997 cosine similarity. However, all existing implementations assume full-precision (FP16/BF16) model weights. We discover that applying TurboQuant to the dominant deployment format---GGUF quantized models (Q4\_K\_M)---produces catastrophic output degradation: language mixing, semantic collapse, and gibberish within 50 tokens. Per-key quality metrics (cos\_sim=0.996) fail to detect this failure because compound quantization error propagates \emph{multiplicatively} through softmax attention, not additively through reconstruction.

We present three contributions: (1)~A \textbf{3-Fix framework} comprising sink token preservation, Past-Only Quantization, and cache state management that eliminates compound error without retraining---producing the first working TurboQuant on GGUF models (300+ token coherent multilingual output). (2)~\textbf{SRHT-based QJL} replacing the paper's $O(d^2)$ dense random projection with $O(d \log d)$ structured projection, yielding 115$\times$ speedup and +4.5~dB SNR improvement. (3)~\textbf{Adaptive QJL}---context-length-aware error correction that reduces attention KL divergence by 2.9$\times$, contradicting the community consensus that ``QJL always hurts.''

\noindent\textbf{Code:} \url{https://github.com/onur-gokyildiz-bhi/tq-kv} \quad \textbf{Crate:} \texttt{crates.io/crates/tq-kv}~v0.4.0
\end{abstract}

%=============================================================================
\section{Introduction}
%=============================================================================

\subsection{The Memory Wall in LLM Inference}

The structural evolution of large language models (LLMs) has reached a critical juncture where the primary impediment to further capability is the physical limitations of the hardware memory hierarchy~\citep{google2026turboquant}. As model parameters have scaled to hundreds of billions, the attention mechanism has created a secondary memory crisis: the Key-Value (KV) cache. During autoregressive decoding, the model stores hidden states of all previous tokens, consuming massive amounts of High Bandwidth Memory. For a 70B model at 32K context, the KV cache alone exceeds 20~GB---half the capacity of most consumer GPUs.

\subsection{The Gap Between Theory and Practice}

To address weight memory, the industry adopted GGUF and other post-training quantization (PTQ) schemes that compress FP16 weights to 4--5 bits~\citep{gguf}. TurboQuant~\citep{zandieh2026turboquant} addresses the KV cache side, compressing keys to 2--4 bits using randomized rotations and optimal scalar quantization. However, a significant gap exists: \textbf{TurboQuant's theoretical guarantees assume FP16 inputs, while production models are served in quantized formats.}

When TurboQuant is applied to the output of already-quantized weights, \emph{compound quantization error} occurs---a phenomenon where two stages of quantization noise interact multiplicatively through softmax, producing catastrophic output degradation~\citep{bondarenko2023softmax, wkvquant}.

\textbf{Concurrent work:} TheTom~\citep{theTom2026turbo4} independently identified and fixed implementation bugs in turbo4 on Apple Silicon (Metal), achieving PPL~6.125 on FP16 models. Our work is complementary: we address the distinct problem of compound error when TurboQuant is applied to \emph{already-quantized} GGUF models---a failure mode that persists even with correct TurboQuant implementation.

\subsection{Contributions}

\begin{enumerate}
    \item \textbf{First identification and solution} of compound quantization error in TurboQuant on GGUF models, via a 3-Fix framework requiring no retraining (Section~\ref{sec:3fix}).
    \item \textbf{SRHT-based QJL} achieving $O(d \log d)$ complexity with 115$\times$ speedup and superior quality over the paper's dense projection (Section~\ref{sec:srht}).
    \item \textbf{Adaptive QJL}---context-aware error correction with empirical evidence that structured QJL improves attention accuracy at all tested context lengths (Section~\ref{sec:adaptive}).
    \item \textbf{tq-kv}---a Pure Rust implementation (10K~LOC, zero C/C++ dependencies) validated on five model architectures on consumer hardware (Section~\ref{sec:impl}).
\end{enumerate}

%=============================================================================
\section{Background}
%=============================================================================

\subsection{TurboQuant: PolarQuant and Lloyd-Max Quantization}

TurboQuant transforms vector quantization into scalar quantization through a two-stage process rooted in Shannon's source coding theory~\citep{zandieh2026turboquant}. For a key vector $\mathbf{k} \in \mathbb{R}^d$:

\textbf{Stage 1---Random Rotation:}
\begin{equation}
    \mathbf{x} = \frac{1}{\sqrt{d}} \mathbf{H} \mathbf{D} \mathbf{k}
\end{equation}
where $\mathbf{H}$ is the normalized Walsh-Hadamard matrix and $\mathbf{D} = \mathrm{diag}(s_1, \ldots, s_d)$ with $s_i \in \{+1, -1\}$ from a shared seed. After rotation, coordinates follow $\mathcal{N}(0, \sigma^2)$.

\textbf{Stage 2---Lloyd-Max Scalar Quantization:}
Each coordinate is independently quantized using pre-computed optimal centroids minimizing MSE for the Gaussian distribution, with $2^b$ centroids for $b$-bit quantization.

\textbf{Our contribution---Adaptive Sigma:} The paper uses fixed $\sigma = 1/\sqrt{d}$. We use per-vector:
\begin{equation}
    \sigma_i = \|\mathbf{k}_i\| / \sqrt{d}
\end{equation}
scaling the codebook to each vector's actual variance.

\subsection{GGUF Weight Quantization}

GGUF Q4\_K\_M uses mixed-precision block quantization with 6-bit scaling factors per mini-block, achieving $\sim$4.5 effective bits~\citep{gguf}. The key projection produces:
\begin{equation}
    \mathbf{k} = \mathbf{W}_{K}^{q} \mathbf{x} + \boldsymbol{\epsilon}_w
\end{equation}
When this noisy key enters TurboQuant, a second noise stage is added: $\mathbf{k}_{\text{tq}} = \text{TQ}(\mathbf{k}) + \boldsymbol{\epsilon}_{\text{tq}}$.

%=============================================================================
\section{The Compound Quantization Error Problem}
\label{sec:compound}
%=============================================================================

\subsection{Failure Mode}

We apply TurboQuant 4-bit to Qwen~2.5 7B in GGUF Q4\_K\_M format. Without TurboQuant, coherent multilingual output. With TurboQuant---even compressing just 2 of 28 layers---catastrophic degradation:

\begin{quote}
\small
\textbf{Without TQ:} ``Merhaba! Bug\"un\"un hava durumu hakk{\i}nda bilgi vermek i\c{c}in l\"utfen hangi \c{s}ehrinizi belirtir misiniz?'' \\[0.5em]
\textbf{With TQ (before fix):} ``Selam! Iyiyim, size Nasil yardimci.GetAxis pestic? kukka Rencontre\ldots'' \emph{[degenerates into 5+ languages within 50 tokens]}
\end{quote}

\textbf{Crucially, per-key quality appears high:} cosine similarity 0.996, SNR 20.3~dB.

\subsection{Root Cause: Softmax Bias Amplification}

The attention score for token $i$ is:
\begin{equation}
    \text{attn}_i = \frac{\exp(\mathbf{q} \cdot \mathbf{k}_i / \sqrt{d})}{\sum_j \exp(\mathbf{q} \cdot \mathbf{k}_j / \sqrt{d})}
\end{equation}

\citet{bondarenko2023softmax} demonstrated that softmax is \textbf{8$\times$ more sensitive} to quantization noise than other activations. The noise creates systematic (not random) bias through the exponential. With 2000+ compressed keys across 24 layers, accumulated multiplicative bias shifts the attention distribution until semantic coherence collapses.

\subsection{Industry Confirmation}

\begin{table}[h]
\centering
\caption{Documented failures of KV cache quantization on quantized-weight models.}
\label{tab:failures}
\begin{tabular}{@{}llll@{}}
\toprule
System & Model & KV Format & Failure \\
\midrule
vLLM \#10411 & Qwen 2.5 7B Q5\_K\_M & FP8 KV & Infinite repetition \\
llama.cpp \#10697 & Llama 3.3 70B Q4\_K\_M & q8\_0 KV & Incorrect output \\
ik\_llama.cpp \#1142 & Various & q4/q5 KV & Garbled text \\
gguf-runner~\citep{jens2026ggufrunner} & Qwen3 30B Q4\_K\_M & $\sim$3-bit KV & No quality metrics \\
All other TurboQuant impl. & Various & 4-bit KV & FP16 only tested \\
\bottomrule
\end{tabular}
\end{table}

%=============================================================================
\section{The 3-Fix Framework}
\label{sec:3fix}
%=============================================================================

Three inference-time corrections requiring no retraining, calibration, or architectural changes.

\subsection{Fix 1: Sink Token Preservation}

Attention sink tokens receive disproportionate attention weight regardless of semantic relevance~\citep{xiao2024sinks}. KVSink~\citep{kvsink} shows preserving their precision reduces error by up to 81\%.

First $N$ tokens' keys are stored in a separate FP16 buffer:
\begin{equation}
    \mathbf{K}_{\text{attn}} = [\underbrace{\mathbf{K}_{\text{sink}}}_{\text{FP16}} \;|\; \underbrace{\mathbf{K}_{\text{compressed}}}_{\text{decompressed}} \;|\; \underbrace{\mathbf{K}_{\text{current}}}_{\text{FP16}}]
\end{equation}

Default $N{=}4$. Overhead: $N \times n_{\text{kv}} \times d \times 4$ bytes $\approx$ 8~KB per layer.

\subsection{Fix 2: Past-Only Quantization (POQ)}

Adapted from WKVQuant~\citep{wkvquant}: the current token's key remains FP16 during attention. It is compressed for future tokens, but attention at time $t$ uses the lossless original:

At time $t$: $[\text{sink}_{\text{FP16}} \;|\; \text{past}_{\text{compressed}} \;|\; \text{current}_t^{\text{FP16}}]$

At time $t{+}1$: current$_t$ moves to compressed cache; current$_{t+1}$ is FP16.

Zero memory overhead---the current token is already resident in registers.

\subsection{Fix 3: Cache State Management}

A previously unidentified bug: \texttt{clear\_cache()} reset the position counter but not compressed KV state. Subsequent conversations appended to stale cache, mixing keys from unrelated contexts with mismatched RoPE positions. Fix: hard reset when \texttt{index\_pos == 0}.

\subsection{Results}

\begin{table}[h]
\centering
\caption{Cumulative impact of 3-Fix on Qwen~2.5 7B Q4\_K\_M with 4-bit TurboQuant.}
\label{tab:3fix}
\begin{tabular}{@{}lcc@{}}
\toprule
Configuration & Max Coherent Tokens & Language Mixing \\
\midrule
No TurboQuant (baseline) & unlimited & none \\
TurboQuant (no fixes) & $\sim$50 & severe (5+ languages) \\
+ Cache reset only & $\sim$100 & moderate \\
+ Cache reset + POQ & $\sim$200 & mild \\
\textbf{+ All 3 fixes (3-Fix)} & \textbf{300+} & \textbf{none} \\
\bottomrule
\end{tabular}
\end{table}

\textbf{With all three fixes, tq-kv produces 300+ token coherent multilingual output on GGUF Q4\_K\_M---the first working TurboQuant on quantized-weight models.}

%=============================================================================
\section{SRHT-Based QJL}
\label{sec:srht}
%=============================================================================

\subsection{The QJL Bottleneck}

TurboQuant's optional QJL error correction projects the quantization residual through a dense Rademacher matrix at $O(d^2)$ cost. Our benchmarks show 29$\times$ compress and 128$\times$ decompress slowdown, making QJL impractical.

\subsection{SRHT Replacement}

We replace the dense matrix $\mathbf{R}$ with a Subsampled Randomized Hadamard Transform:
\begin{align}
    \text{Dense:}\quad & \mathbf{p} = \mathbf{R} \boldsymbol{\epsilon} & O(md) \\
    \text{SRHT:}\quad & \mathbf{p} = \sqrt{d/m} \cdot \mathbf{S} \mathbf{H} \mathbf{D} \boldsymbol{\epsilon} & O(d \log d + m)
\end{align}
where $\mathbf{D}$ = random signs, $\mathbf{H}$ = Walsh-Hadamard, $\mathbf{S}$ = row subsample~\citep{ailon2009fast}. All primitives reuse our existing Hadamard module. The batch path shares $\mathbf{D}$ across vectors and reuses buffers, eliminating per-vector allocation.

\subsection{Results}

\begin{table}[h]
\centering
\caption{SRHT vs Dense QJL. 32,768 vectors, $d{=}128$, 4-bit, release build.}
\label{tab:srht}
\begin{tabular}{@{}lccc@{}}
\toprule
Metric & Dense QJL & SRHT QJL & No QJL \\
\midrule
Compress time & 2196~ms & 100~ms & 69~ms \\
Decompress time & 2941~ms & 39~ms & 23~ms \\
Compress overhead & 29$\times$ & \textbf{1.45$\times$} & 1.0$\times$ \\
Decompress overhead & 128$\times$ & \textbf{1.7$\times$} & 1.0$\times$ \\
SNR vs no-QJL & +1.2~dB & \textbf{+4.5~dB} & baseline \\
Cosine Similarity & 0.9966 & \textbf{0.9984} & 0.9956 \\
\bottomrule
\end{tabular}
\end{table}

The +4.5~dB SNR improvement over dense QJL is unexpected. We hypothesize that SRHT's structured projection distributes energy more uniformly, reducing variance in the sign quantization step.

%=============================================================================
\section{Adaptive QJL}
\label{sec:adaptive}
%=============================================================================

\subsection{The ``QJL Hurts'' Consensus}

Five independent groups have confirmed that QJL hurts quality: ikawrakow/TheTom (turboquant\_plus on Metal)~\citep{theTom2026turbo4}, spiritbuun (CUDA), scos-lab, Arclabs001/YATQ, and the paper's own ablation on OpenReview~\citep{llamacpp_tq}. TheTom's analysis is definitive: ``QJL eliminates bias but explodes variance. Softmax amplifies variance. So you're spending 1 bit per value to make your model dumber.'' Their solution---drop QJL, use all 4 bits for 16-centroid PolarQuant---achieved PPL~6.125 (+0.23\% vs q8\_0) on Qwen~3.5 35B MoE~\citep{theTom2026turbo4}.

The key insight, as TheTom articulates: ``The paper optimized for MSE. Attention optimizes through softmax. Different objective, different answer.'' More centroids consistently beat error correction at every tested configuration.

\subsection{Our Counterargument: Dense vs Structured Projection}

These findings all used the paper's \emph{dense Rademacher} QJL, which has high variance from random projection. We observe that SRHT QJL has fundamentally different variance properties due to structured (Hadamard-based) projection. We measure \textbf{attention KL divergence} (not just MSE) on synthetic Gaussian data:

\begin{table}[h]
\centering
\caption{Attention KL divergence: no QJL vs SRHT QJL across context lengths.}
\label{tab:adaptive}
\begin{tabular}{@{}rccc@{}}
\toprule
Context Length & KL(no QJL) & KL(SRHT QJL) & Reduction \\
\midrule
64 & 0.000270 & 0.000093 & 2.9$\times$ \\
256 & 0.000283 & 0.000099 & 2.9$\times$ \\
1,024 & 0.000289 & 0.000103 & 2.8$\times$ \\
4,096 & 0.000288 & 0.000102 & 2.8$\times$ \\
8,192 & 0.000287 & 0.000102 & 2.8$\times$ \\
16,384 & 0.000288 & 0.000103 & 2.8$\times$ \\
\bottomrule
\end{tabular}
\end{table}

On synthetic data, SRHT QJL consistently reduces KL divergence by $\sim$2.9$\times$ at all context lengths. Top-5 token accuracy improves from 80--85\% to 87--95\%.

\textbf{Important caveat:} These results are on synthetic Gaussian data. TheTom's real-model evidence (PPL measurements on Qwen~3.5 35B) shows dense QJL hurting quality. Whether SRHT's lower variance changes this tradeoff on real models remains an open question requiring perplexity validation at extended context lengths. We present these results as preliminary evidence that the projection structure matters, not as a definitive recommendation.

\subsection{Adaptive Activation}

Given the tension between our synthetic results (SRHT QJL helps) and community real-model evidence (dense QJL hurts), we introduce a conservative \texttt{QjlMode::Adaptive\{threshold\}}: QJL off below threshold (following TheTom's recommendation for short context), on above (where accumulated bias may outweigh variance). Default threshold: 4096 tokens. This hedges against both failure modes until real-model validation resolves the question.

%=============================================================================
\section{Implementation}
\label{sec:impl}
%=============================================================================

\subsection{System Architecture}

tq-kv is 10,000 lines of Pure Rust with zero C/C++ dependencies:

\begin{itemize}
    \item \textbf{Library} (crates.io v0.4.0): 2/3/4-bit Lloyd-Max, SRHT QJL, fused attention (AVX2+FMA SIMD), TurboKvCache (candle drop-in), C~FFI
    \item \textbf{Engine}: GenericTurboModel for 5 architectures (Qwen2, Llama, Mistral, Phi3, Gemma2) with GGUF auto-detection and 3-Fix framework
    \item \textbf{Platform}: OpenAI-compatible API (axum), model hub, Web~UI, auto-TQ
\end{itemize}

\subsection{Fused Attention}

Because the Walsh-Hadamard Transform is orthogonal: $\langle \mathbf{q}, \mathbf{R}^{-1} \mathbf{k}_c \rangle = \langle \mathbf{R}\mathbf{q}, \mathbf{k}_c \rangle$. The query is rotated once; scores are computed via centroid table lookup against compressed indices---6$\times$ faster than decompress-then-dot with AVX2+FMA SIMD.

\subsection{VRAM Impact}

\begin{table}[h]
\centering
\caption{VRAM savings with TurboQuant compression (keys only).}
\label{tab:vram}
\begin{tabular}{@{}lrrrr@{}}
\toprule
Model (4K context) & FP16 KV & TQ 4-bit & TQ 2-bit & Max Savings \\
\midrule
Qwen 2.5 7B & 256~MB & 48~MB & 18~MB & 14.2$\times$ \\
Qwen 2.5 72B & 640~MB & 120~MB & 45~MB & 14.2$\times$ \\
Llama 3.1 70B (32K) & 20~GB & 5.3~GB & 1.4~GB & 14.2$\times$ \\
\bottomrule
\end{tabular}
\end{table}

%=============================================================================
\section{Limitations and Future Work}
%=============================================================================

\textbf{Fused kernel divergence:} The fused centroid-lookup path shows slight numerical divergence from the decompress path through softmax. Currently defaulting to decompress; stable fused kernels needed.

\textbf{Per-channel scaling:} SmoothQuant-style per-channel normalization could further reduce compound error by handling outlier channels before rotation.

\textbf{Calibrated codebook:} Lloyd-Max centroids assume Gaussian. Calibrating from Q4 model activations could improve quality, particularly at 2-bit.

\textbf{Learned rotations:} SpinQuant~\citep{spinquant} shows learned rotations outperform random Hadamard by up to 45\%. Integration could enable stable 2-bit on GGUF.

\textbf{Real-model QJL validation:} Our adaptive QJL results are on synthetic data. Extended perplexity benchmarks at 4K--32K context are needed.

%=============================================================================
\section{Conclusion}
%=============================================================================

We demonstrate that TurboQuant KV cache compression fails on GGUF quantized models---the dominant deployment format---due to compound quantization error. This failure affects all existing implementations. Our 3-Fix framework solves this at inference time without retraining, producing the first working TurboQuant on GGUF: 300+ token coherent output where all other implementations produce gibberish.

SRHT-based QJL achieves 115$\times$ speedup with +4.5~dB SNR improvement, and adaptive QJL reduces attention KL divergence by 2.9$\times$. These contributions open a new research direction: \textbf{compound-aware KV cache compression}---methods designed for the quantized models that dominate real-world deployment.

%=============================================================================
\bibliographystyle{plainnat}
\bibliography{references}

\end{document}
